\chapter{Combining Decisions and Loops in Python}

\section{Introduction: Real-Life Hook}

Students often face repeated decisions:

\begin{itemize}
    \item Checking multiple assignments: "Is it done or not?"  
    \item Asking friends for answers: "Yes or No" multiple times  
    \item Counting how many books are unread in a shelf  
\end{itemize}

Python can handle such tasks efficiently by combining \textbf{loops} and \textbf{if/else statements}.

\section{Concept: Decision-Making Inside Loops}

\begin{definition}{Decision Inside Loop}
A program can check conditions repeatedly by placing \textbf{if/else statements inside loops}.  
This allows programs to respond dynamically during repetition.
\end{definition}

\subsection{Example: Checking Multiple Marks}

Ask a user for marks of 5 students and determine pass/fail:

\begin{verbatim}
for i in range(1, 6):
    marks = int(input("Enter marks of student " + str(i) + ": "))
    if marks >= 40:
        print("Student", i, "passed")
    else:
        print("Student", i, "failed")
\end{verbatim}

- Loop repeats 5 times  
- Decision inside the loop determines pass/fail for each student

\section{Example: Guess the Number Game}

\begin{verbatim}
secret = 7
guess = 0
while guess != secret:
    guess = int(input("Guess the number (1-10): "))
    if guess < secret:
        print("Too low, try again!")
    elif guess > secret:
        print("Too high, try again!")
    else:
        print("Congratulations! You guessed it!")
\end{verbatim}

- Loop continues until the correct guess  
- Conditional statements guide the user

\section{Step-by-Step Diagram: Loop with Decisions}

\begin{center}
\begin{tikzpicture}[node distance=1.5cm]
\node (start) [draw, ellipse, fill=green!20] {Start Program};
\node (loop) [draw, diamond, fill=yellow!20, below of=start, yshift=-1cm] {Loop Condition True?};
\node (input) [draw, rectangle, fill=orange!20, below of=loop, yshift=-1cm] {Get Input / Task};
\node (decision) [draw, diamond, fill=yellow!20, below of=input, yshift=-1cm] {Decision Condition True?};
\node (trueblock) [draw, rectangle, fill=orange!20, below left of=decision, xshift=-2cm, yshift=-1cm] {Execute if True};
\node (falseblock) [draw, rectangle, fill=orange!20, below right of=decision, xshift=2cm, yshift=-1cm] {Execute if False};
\node (update) [draw, rectangle, fill=orange!20, right of=loop, xshift=4cm] {Update Counter/Variable};
\node (end) [draw, ellipse, fill=green!20, left of=loop, xshift=-4cm] {End Program};

\draw[->] (start) -- (loop);
\draw[->] (loop.south) -- node[midway, right] {Yes} (input.north);
\draw[->] (input.south) -- (decision.north);
\draw[->] (decision.west) -| node[near start, above] {Yes} (trueblock.north);
\draw[->] (decision.east) -| node[near start, above] {No} (falseblock.north);

% Connect blocks back to update
\draw[->] (trueblock.south) |- (0,-9) -| (update.south);
\draw[->] (falseblock.south) |- (0,-9) -| (update.south);

% Update back to loop
\draw[->] (update.west) -- (loop.east);

% End program
\draw[->] (loop.west) -- node[midway, above] {No} (end.east);

\end{tikzpicture}
\end{center}

\section{Mini-Project: Simple Quiz Game}

Create a program that asks 3 questions and keeps score:

\begin{verbatim}
score = 0

q1 = input("What is 2 + 2? ")
if q1 == "4":
    print("Correct!")
    score += 1
else:
    print("Wrong!")

q2 = input("What is 5 * 3? ")
if q2 == "15":
    print("Correct!")
    score += 1
else:
    print("Wrong!")

q3 = input("What is 10 - 7? ")
if q3 == "3":
    print("Correct!")
    score += 1
else:
    print("Wrong!")

print("Your total score is:", score)
\end{verbatim}

\section{Key Concepts}

\begin{itemize}
    \item Loops can repeat tasks multiple times.  
    \item Decisions inside loops make programs interactive and responsive.  
    \item Combining loops and if/else reduces repetitive code.  
    \item Variables can store state (like score) across iterations.  
\end{itemize}

\section{End-of-Chapter Exercises}

\begin{exercise}{Very Short Questions}
\begin{enumerate}
    \item What is the purpose of placing if/else inside a loop?  
    \item Can loops run without decisions?  
    \item Give an example of a repeated decision in real life.  
\end{enumerate}
\end{exercise}

\begin{exercise}{Short Answer Questions}
\begin{enumerate}
    \item Write a Python program that asks 5 users for their ages and prints whether each can vote.  
    \item Explain how the guess-the-number game works.  
    \item Why do we update variables inside loops?  
\end{enumerate}
\end{exercise}

\begin{exercise}{Activity-Based Questions}
\begin{enumerate}
    \item Create a simple arithmetic quiz with 3 questions and calculate the score using loops and if/else.  
    \item Modify the guess-the-number game to give a hint about whether the guess is even or odd.  
    \item Write a program that asks for 5 exam marks and prints “Excellent” for marks $\ge$ 90, “Good” for 70–89, else “Needs Improvement”.  
\end{enumerate}
\end{exercise}

\section{Teacher Notes}

\marginnote{
\textbf{Learning Objectives:}  
\begin{itemize}
    \item Students combine loops and if/else statements effectively.  
    \item Students understand dynamic, interactive programs.  
    \item Students can create games and quizzes using loops and decisions.  
\end{itemize}

\textbf{Teaching Tips:}  
\begin{itemize}
    \item Demonstrate examples interactively in class.  
    \item Encourage students to test different inputs to explore all branches.  
\end{itemize}

\textbf{Assessment:}  
\begin{itemize}
    \item Observe understanding of loop-decision combination.  
    \item Check correctness of mini-projects like quiz and number-guessing games.  
\end{itemize}
}